\subsection{Binary-Polyhedral Maximality and the $SU(2)$ Thread as Structural Attractor}
  \label{subsec:binary-polyhedral-maximality}

  The classification obtained in the previous section identifies a restricted
  set of finite groups capable of supporting admissible neutral generating
  sets.
  Among these structures, the binary polyhedral groups occupy a distinguished
  position.

  While the quaternion and binary dihedral families arise from anisotropic
  configurations of the second-moment tensor, the binary polyhedral groups
  $2O$ and $2I$ correspond to isotropic configurations in which
  \[
    M \propto I .
  \]
  In such configurations the generators are distributed uniformly on the
  sphere and the admissibility constraints are satisfied in the most symmetric
  possible manner.

  This enhanced symmetry has important consequences for the admissible
  spectral hierarchy.
  Isotropic configurations maximise the uniformity of the Gram relations and
  therefore minimise the structural tension between different representation
  sectors.

  As a result, the binary polyhedral groups realise particularly stable
  arrangements of admissible generators.
  In comparison with groups whose generator configurations are anisotropic,
  their admissible spectral sectors remain more robust under variations of the
  generating set.

  This property becomes even clearer when one considers global measures of
  admissibility such as the penalised spectral capacity $C_\alpha$.
  When constructive alignment between sectors is penalised, binary polyhedral
  groups systematically dominate their non-binary counterparts obtained by
  quotienting the central $\mathbb{Z}_2$.

  In this sense the double-cover structure characteristic of binary groups is
  not merely a representational convenience.
  Rather, it reflects a deeper compatibility between the admissibility
  constraints and the spinorial geometry of the underlying representations.

  The binary polyhedral groups therefore form natural attractors within the
  space of finite relational structures compatible with admissible spectral
  dynamics.

  Moreover, these groups are embedded in the continuous group $SU(2)$.
  This observation suggests that the structures identified above are not
  isolated algebraic curiosities but rather finite manifestations of a deeper
  spinorial framework.

  The role of $SU(2)$ as a structural limit of these constructions will be
  examined in the following section, where sequences of increasingly refined
  relational graphs are considered.
